\documentclass[10pt,a4paper]{article}

% ----------------- Paquetes -------------------%
\usepackage[T1]{fontenc}
\usepackage[utf8]{inputenc}
\usepackage[a4paper,margin=2.5cm]{geometry}
\usepackage{mathptmx}             
\usepackage{changepage} 
\usepackage{setspace}
\usepackage[spanish]{babel}
\usepackage[labelfont=bf,labelsep=space]{caption}
\usepackage{amssymb}
\usepackage{amsmath}
\usepackage{ebgaramond}
\usepackage{placeins}
\usepackage{etoolbox}
\usepackage{setspace}
\usepackage{parskip} 
\usepackage{tcolorbox}
\usepackage{needspace}
\usepackage{xparse}
\usepackage{ocgx2}
\usepackage{hyperref}
\usepackage{hypcap}
\usepackage{enumitem}


%%%%%%%%%%%%%%%%%%%%%%%%%%%%%%%%%%%%%%%%%%%%%%%%%%%%%%%%%%%%%%%%
% --- Fija primero tus valores globales "buenos" ---
\setstretch{1.2}                 % Interlineado global
\setlength{\parskip}{0.7\baselineskip}
\setlength{\parindent}{0pt}
\newlength{\GlobalParskip}
\setlength{\GlobalParskip}{\parskip}

\newcounter{eqnum}[subsection]
\renewcommand{\theeqnum}{\arabic{eqnum}}
\newcounter{alphaitem}
\newcounter{numitem}

%% ------------------- Recuadros----------------------%%
\newcommand{\puntosclave}[1]{%
  \begin{tcolorbox}[
    colframe=black,
    title={Puntos clave},
    before upper={\setlength{\parindent}{0.5cm}\setlength{\parskip}{0.5\baselineskip}}
  ]
  #1
  \end{tcolorbox}
}

\newcommand{\modificaciones}[1]{
  \begin{tcolorbox}[
    colback=white,
    colframe=red,
    title={Modificaciones a realizar},
    before upper={\setlength{\parindent}{0.5cm}\setlength{\parskip}{0.5\baselineskip}}
  ]
  #1
  \end{tcolorbox}
}

%% ------------------- Preguntas TEST----------------------%%
% Opciones: radio (single-choice)
\NewDocumentCommand{\OptRadio}{m m m}{%
  \noindent\CheckBox[radio,name=#1,value=#2,width=1.2ex,height=1.2ex]{}%
  \;\textbf{#2.}~#3\par
}

% Opciones: checkbox (multi-choice)
\NewDocumentCommand{\OptCheck}{m m}{%
  \noindent\CheckBox[name=#1,width=1.2ex,height=1.2ex,bordercolor={0 0 0}]{}%
  \;#2\par
}

% Pregunta single-choice (radio)
% Args: 1=id, 2=enunciado, 3=opciones, 4=solución+justificación, 5=imagen (o {}), 6=origen
\NewDocumentCommand{\PreguntaTestSingle}{m m m m m m}{%
  \par\medskip
  \noindent\textbf{#1.}~#2\par\smallskip

    % Imagen (si no es {})
    \IfBlankTF{#5}{}{%
      \begin{center}
        \includegraphics[width=0.75\linewidth]{#5}
      \end{center}
    }

    \begin{Form}
      #3
    \end{Form}

    \par\smallskip
    \switchocg{sol-#1}{\fbox{Mostrar / ocultar solución}}%
    \hfill{\footnotesize #6}\par\smallskip

    \begin{ocg}{Solución}{sol-#1}{0}
      \noindent #4
    \end{ocg}
  \par\medskip
}

% Pregunta multi-choice (checkboxes)
\NewDocumentCommand{\PreguntaTestMulti}{m m m m m m}{%
  \par\medskip
  \noindent\textbf{#1.}~#2\par\smallskip

    % Imagen (si no es {})
    \IfBlankTF{#5}{}{%
      \begin{center}
        \includegraphics[width=0.75\linewidth]{#5}
      \end{center}
    }
    \begin{Form}
      #3
    \end{Form}

    \par\smallskip
    \switchocg{sol-#1}{\fbox{Mostrar / ocultar solución}}%
    \hfill{\footnotesize #6}\par\smallskip

    \begin{ocg}{Solución}{sol-#1}{0}
      \noindent #4
    \end{ocg}
  \par\medskip
}

%% ------------------- Listas----------------------%%
    % --- Lista alfabética ---
    \newenvironment{la}
      {\begin{list}{\alph{alphaitem})}{%
          \usecounter{alphaitem}%
          \setlength{\leftmargin}{1cm}%
          \setlength{\parskip}{3pt}% 
          \setlength{\itemsep}{0pt}% ← espacio entre ítems
          \setlength{\parsep}{0pt}%  ← párrafos dentro de un ítem
      }}
      {\end{list}}
      
    % --- Lista numérica ---
    \newenvironment{lnum}
      {\begin{list}{\arabic{numitem}.}{%
          \usecounter{numitem}%
          \setlength{\leftmargin}{1cm}%
          \setlength{\parskip}{3pt}% 
          \setlength{\itemsep}{0pt}% ← espacio entre ítems
          \setlength{\parsep}{0pt}%  ← párrafos dentro de un ítem
      }}
      {\end{list}}
      
% ------------------- Imágenes----------------------%
    \graphicspath{{Img/}}% Rutas por defecto 
    % --- Imagen adaptativa ---
    % Registros de dimensión definidos una sola vez
    \newdimen\imgcap
    \newdimen\imgmin
    \newdimen\imgmax
    \newdimen\imgremain
    \newdimen\imgh
    \newdimen\imgneed
    
    \newcommand{\imagenfit}[3]{%
      \addvspace{10pt}% separación antes
      \par\begingroup
        % Parámetros de composición (asignaciones, no \newdimen)
        \imgcap=1\baselineskip            % alto estimado de título+fuente
        \imgmin=0.15\textheight
        \imgmax=0.15\textheight
        \imgremain=\pagegoal
        \ifdim\pagegoal=\maxdimen \imgremain=0pt\fi
        \advance\imgremain by -\pagetotal
        \advance\imgremain by -\imgcap
        % Decisión de altura destino
        \ifdim\imgremain<\imgmin
          \newpage
          \imgh=0.20\textheight
        \else
          \imgh=\imgremain
          \ifdim\imgh>\imgmax \imgh=\imgmax \fi
        \fi
        % Reservar espacio para evitar cortes del bloque completo
        \imgneed=\imgh \advance\imgneed by \imgcap
        \Needspace{\imgneed}
        % Cuerpo
        \noindent\begin{minipage}{\textwidth}
          \centering
          \captionof{figure}{\textemdash\ #2}\par\vspace{0.3em}% título arriba
          \includegraphics[width=\linewidth,height=\imgh,keepaspectratio]{\detokenize{#1}}\par
          \vspace{0.3em}\small\textit{Fuente: }#3% fuente abajo
        \end{minipage}%
      \endgroup
      \par\addvspace{10pt}%
    }

%------------ Bloque de RECUADRO ESPECIAL -------------%
    \newenvironment{specialblock}[1]{%
      \begin{quote} % crea sangría a izquierda y derecha
      \small % texto más pequeño
      \noindent\textbf{#1}\\[-0.3em] % título en negrita
      \rule{\linewidth}{0.4pt}\\[0.6em]}{
      \end{quote}}
      
%-------------------------- Portada -----------------------%
\newcommand{\oldtitle}[1]{\def\OldTitle{#1}}
\newcommand{\changerationale}[1]{\def\ChangeWhy{#1}}

\newcommand{\changebox}{%
\begin{tcolorbox}[title={Historial de cambios del tema}]
\textbf{Título anterior:} \OldTitle\par
\textbf{Motivo del cambio:} \ChangeWhy
\end{tcolorbox}
}

%-------------------------- Author Font -----------------------%
    \newcommand{\authorfont}[1]{{\fontfamily{EBGaramond-TLF}\selectfont #1}}

%-------------------------- Anexos -----------------------%
    \makeatletter
    \newcounter{anexo}
    \renewcommand{\theanexo}{\Roman{anexo}}
    \newcommand{\anexo}[1]{%
      \clearpage
      \phantomsection
      \refstepcounter{anexo}%
      \noindent\textbf{Anexo \theanexo.- }#1\par\medskip
      \addcontentsline{stc}{section}{Anexo \theanexo.- #1}%
    }
    \makeatother

    %-------------------------- Lista de figuras (personal) -----------------------%
\makeatletter
\newcommand{\MyListOfFigures}{%
  \clearpage
  \phantomsection
  {\centering\bfseries\Large Índice de figuras\par}%
  \medskip
  % Entrada en nuestro índice de contenido (stc)
  \addcontentsline{stc}{section}{Índice de figuras}%
  % Recuperación del listado (sin cabecera propia)
  \begingroup
    \parindent=0pt\parskip=2pt
    \@starttoc{lof}%
  \endgroup
}
\makeatother

%-------------------------- Bibliografía -----------------------%
\makeatletter
% Contador para las referencias
\newcounter{myref}
% Cabecera de la bibliografía, entra en el índice 'stc'
\newcommand{\MyBibliography}{%
  \clearpage
  \phantomsection
  {\centering\bfseries\Large Bibliografía y referencias\par}%
  \medskip
  \addcontentsline{stc}{section}{Bibliografía y referencias}%
  \setcounter{myref}{0}%
}
\newcommand{\MyBibItem}[4]{%
  \refstepcounter{myref}%
  \noindent[\themyref]\ #1.\ \emph{#2}. #3. #4\par
}
\makeatother

%--------- Establecer cuatro niveles de índice --------%
    \setcounter{tocdepth}{4}   
    \setcounter{secnumdepth}{4}% numera hasta \paragraph

%%%%%%%%%%%%%%%%%%%%%%%%%%%%%%%%

\makeatletter

    %--------- Interlineado Global --------%
    \edef\GlobalBaseLineStretch{\baselinestretch}
    
    %--------- Espaciado de seccinones --------%
    \renewcommand\section{\@startsection{section}
      {1}{\z@}
      {2.5ex \@plus 1ex \@minus .2ex} % espacio antes
      {1.5ex \@plus .2ex}             % espacio después
      {\normalfont\Large\bfseries}    % formato del título
    }
    \renewcommand\subsection{\@startsection{subsection}{2}{\z@}
      {3ex \@plus 0.8ex \@minus .2ex}
      {1.5ex \@plus .2ex}
      {\normalfont\large\bfseries}
    }
    \renewcommand\subsubsection{\@startsection{subsubsection}{3}{\z@}
      {3ex \@plus 0.5ex \@minus .2ex}
      {1ex \@plus .2ex}
      {\normalfont\normalsize\bfseries}
    }
    \renewcommand\paragraph{\@startsection{paragraph}{4}{\z@}%
    {-2ex \@plus -0.5ex \@minus -.2ex}% espacio antes
    {0.8ex \@plus .2ex}% espacio después
    {\normalfont\normalsize\itshape}}% ← cursiva en lugar de negrita

    %--------- Ecuaciones --------%   
    % Variables globales para almacenar títulos y notas
    \gdef\leq@current@section{}%
    \gdef\leq@current@subsection{}%
    \gdef\leq@last@printed@section{}%
    \gdef\leq@last@printed@subsection{}%
    
    % Comando para añadir nota (invisible en el cuerpo)
    \newcommand{\eqnote}[1]{%
      \addcontentsline{leq}{leqnote}{#1}%
    }
    
    % Comando para ecuaciones
    \newcommand{\eqblock}[2]{%
      % Verificar si necesitamos imprimir el título de sección
      \ifx\leq@current@section\leq@last@printed@section\else
        \addcontentsline{leq}{leqsection}{\leq@current@section}%
        \global\let\leq@last@printed@section\leq@current@section
        \gdef\leq@last@printed@subsection{}% Reset subsección al cambiar sección
      \fi
      % Verificar si necesitamos imprimir el título de subsección
      \ifx\leq@current@subsection\leq@last@printed@subsection\else
        \ifx\leq@current@subsection\@empty\else
          \addcontentsline{leq}{leqsubsection}{\leq@current@subsection}%
          \global\let\leq@last@printed@subsection\leq@current@subsection
        \fi
      \fi
      % Ecuación
      \refstepcounter{eqnum}%
      \begin{equation*}
        #1\tag{E.\theeqnum}
      \end{equation*}%
      \addcontentsline{leq}{leqt}{[E.\theeqnum] -- #2}%
      \addcontentsline{leq}{leqm}{#1}%
    }
    
    %%% ----- Índice de ecuaciones ----------%%%
    % Tamaño para []
    \newcommand{\leqIndexFont}{\normalsize}
    \newcommand{\leqTitleFont}{\tiny}
    \newcommand{\leqNoteFont}{\tiny}
    % Cabecera de sección 
    \newcommand*\l@leqsection[2]{%
      \addvspace{25pt}% Mayor separación antes de nueva sección
      \noindent\leqIndexFont
      \makebox[\linewidth]{\rule{.38\linewidth}{0.4pt}\ \ \textbf{  \MakeUppercase{#1}}\ \ \rule{.38\linewidth}{0.4pt}}%
      \par\vspace{-10pt}
    }
    % Cabecera de subsección 
    \newcommand*\l@leqsubsection[2]{%
      \par\addvspace{15pt}%
      \noindent\leqIndexFont #1
      \par\vspace{-7pt}
    }
    % Título de ecuación 
    \newcommand*\l@leqt[2]{%
    \par\addvspace{10pt}%
      \par\noindent\leqTitleFont #1\par
    }
    % Miniatura de ecuación
    \newcommand*\l@leqm[2]{%
      \vspace{0.3\baselineskip}%
      \par\scriptsize\noindent\hspace*{1.5em}\(#1\)\par%
    }
    % Nota de ecuación 
    \newcommand*\l@leqnote[2]{%
      \vspace{0.2\baselineskip}%
      \par\noindent\leqNoteFont\hspace*{1.5em}\textbf{Nota.--} #1\par\vspace{0.5\baselineskip}%
    }
    % Generación del índice
    \newcommand{\mylistofequations}{%
      \clearpage
      \phantomsection
      % Título visual de la lista (ajústalo a tu gusto):
      {\centering\bfseries\Large Índice de ecuaciones\par}
      % Entrada en nuestro índice 'stc':
      \addcontentsline{stc}{section}{Índice de ecuaciones}%
      % Llamada a tu lista real (usa el nombre exacto de tu macro):
      \@starttoc{leq}%
    }
    
    %--------- Captura de títulos de sección --------%
    \let\leq@original@section\section
    \let\leq@original@subsection\subsection
    % Flag para saber si estamos en una sección asterisco
    \newif\ifLeqStarSection
    \LeqStarSectionfalse
    
    % Redefinir \section CON CONTROL DE ASTERISCO
    \renewcommand{\section}{%
      \@ifstar{\leq@section@star}{\@dblarg\leq@section@nostar}%
    }
    \def\leq@section@star#1{%
      \global\LeqStarSectiontrue
      \gdef\leq@current@section{#1}%
      \gdef\leq@current@subsection{}%
      \leq@original@section*{#1}%
      \global\LeqStarSectionfalse
    }
    \def\leq@section@nostar[#1]#2{%
      \global\LeqStarSectionfalse
      \gdef\leq@current@section{#2}%
      \gdef\leq@current@subsection{}%
      \leq@original@section[#1]{#2}%
    }
    % Redefinir \subsection
    \renewcommand{\subsection}{%
      \@ifstar{\leq@subsection@star}{\@dblarg\leq@subsection@nostar}%
    }
    \def\leq@subsection@star#1{%
      \global\LeqStarSectiontrue
      \gdef\leq@current@subsection{#1}%
      \leq@original@subsection*{#1}%
      \global\LeqStarSectionfalse
    }
    \def\leq@subsection@nostar[#1]#2{%
      \global\LeqStarSectionfalse
      \gdef\leq@current@subsection{#2}%
      \leq@original@subsection[#1]{#2}%
    }

    % ----- Restauración de valores Globales de estilo ----------%
    % Flags
    \newif\ifInFootnote
    \InFootnotefalse
    \pretocmd{\@footnotetext}{\InFootnotetrue}{}{}
    \apptocmd{\@footnotetext}{\InFootnotefalse}{}{}
    % Profundidad de listas (soporta anidadas)
    \newcount\MyListDepth \MyListDepth=0
    \AtBeginEnvironment{la}{\advance\MyListDepth by 1}
    \AfterEndEnvironment{la}{\advance\MyListDepth by -1}
    \AtBeginEnvironment{lnum}{\advance\MyListDepth by 1}
    \AfterEndEnvironment{lnum}{\advance\MyListDepth by -1}
    \AtBeginEnvironment{itemize}{\advance\MyListDepth by 1}
    \AfterEndEnvironment{itemize}{\advance\MyListDepth by -1}
    \AtBeginEnvironment{enumerate}{\advance\MyListDepth by 1}
    \AfterEndEnvironment{enumerate}{\advance\MyListDepth by -1}
    % Restauración diferida: hazlo al INICIO del próximo párrafo real
    \newif\if@pendingRestore
    \@pendingRestorefalse
    \newcommand{\RequestRestore}{\global\@pendingRestoretrue}
    \everypar{%
      \if@pendingRestore
        \global\@pendingRestorefalse
        \ifInFootnote\else
          \ifnum\MyListDepth=0
            \setlength{\parskip}{\GlobalParskip}%
            \setstretch{\GlobalBaseLineStretch}%
          \fi
        \fi
      \fi
    }
    % Al final de bloques:
    \AfterEndEnvironment{equation}{\RequestRestore}
    \AfterEndEnvironment{equation*}{\RequestRestore}
    \AfterEndEnvironment{la}{\RequestRestore}
    \AfterEndEnvironment{lnum}{\RequestRestore}
    % Al final de macros:
    \apptocmd{\eqblock}{\RequestRestore}{}{}
    \apptocmd{\imagenfit}{\RequestRestore}{}{}
    
\makeatother

% ===================== ToC propio con 4 niveles (único bloque) =====================
\makeatletter

% --- Escritores: añadimos una línea al .stc en cada cabecera numerada ---
% Guardamos las versiones actuales para llamar al comportamiento normal
\let\MyTOC@orig@section\section
\let\MyTOC@orig@subsection\subsection
\let\MyTOC@orig@subsubsection\subsubsection
\let\MyTOC@orig@paragraph\paragraph

% SECTION
\renewcommand{\section}{%
  \@ifstar{\MyTOC@section@star}{\@dblarg\MyTOC@section@nostar}%
}
\def\MyTOC@section@star#1{%
  \MyTOC@orig@section*{#1}% sin entrada al .stc
}
\def\MyTOC@section@nostar[#1]#2{%
  \MyTOC@orig@section[{#1}]{#2}% comportamiento normal (escribe al .toc)
  % Escribimos también nuestra línea al .stc (texto con \numberline)
  \addcontentsline{stc}{section}{\protect\numberline{\thesection}#2}%
}

% SUBSECTION
\renewcommand{\subsection}{%
  \@ifstar{\MyTOC@subsection@star}{\@dblarg\MyTOC@subsection@nostar}%
}
\def\MyTOC@subsection@star#1{%
  \MyTOC@orig@subsection*{#1}%
}
\def\MyTOC@subsection@nostar[#1]#2{%
  \MyTOC@orig@subsection[{#1}]{#2}%
  \addcontentsline{stc}{subsection}{\protect\numberline{\thesubsection}#2}%
}

% SUBSUBSECTION
\renewcommand{\subsubsection}{%
  \@ifstar{\MyTOC@subsubsection@star}{\@dblarg\MyTOC@subsubsection@nostar}%
}
\def\MyTOC@subsubsection@star#1{%
  \MyTOC@orig@subsubsection*{#1}%
}
\def\MyTOC@subsubsection@nostar[#1]#2{%
  \MyTOC@orig@subsubsection[{#1}]{#2}%
  \addcontentsline{stc}{subsubsection}{\protect\numberline{\thesubsubsection}#2}%
}

% PARAGRAPH
\renewcommand{\paragraph}{%
  \@ifstar{\MyTOC@paragraph@star}{\@dblarg\MyTOC@paragraph@nostar}%
}
\def\MyTOC@paragraph@star#1{%
  \MyTOC@orig@paragraph*{#1}%
}
\def\MyTOC@paragraph@nostar[#1]#2{%
  \MyTOC@orig@paragraph[{#1}]{#2}%
  % Si no numeras los \paragraph, cambia \theparagraph por texto plano sin \numberline
  \addcontentsline{stc}{paragraph}{\protect\numberline{\theparagraph}#2}%
}

% --- Impresor del índice propio (lee el .stc) ---
\newcommand{\MyTOC}{%
  \par\bigskip
  {\centering\bfseries\Large Índice de contenido\par}%
  \medskip
  \begingroup
    \parindent=0pt\parskip=2pt
    \@starttoc{stc}%
  \endgroup
}

\makeatother
% ===================== fin ToC propio =====================


%%%%%%%%%%%%%%


% --- Datos del tema ---
\title{Tema 4.B.34 -- La OMC (I)\\
\large Acuerdos distintos de los de mercancías}
\author{Marcelo Ortega}
\date{Marzo 2026}
\newcommand{\TemaCode}{3B34}

\oldtitle{
    B.35. La OMC. Los acuerdos distintos de los de mercancías
}
\changerationale{
    Sin cambios.
}

\begin{document}


\maketitle
\changebox

\newpage

% --- Página de resumen y modificaciones ---
\puntosclave{ %% Escribir puntos clave Apuntes CECO
     Recomendaciones para actualizar el tema: Como en todos los temas de la oposición, se recomienda profundizar en cada nueva vuelta al temario. Para ello, lo más práctico es buscar temas concretos en la web de la OMC (así como para conocer los resultados de la última Conferencia Ministerial). A veces la revista ICE publica algún contenido de interés. Como consejo personal del autor, busca adaptar el contenido del tema a tus habilidades y fortalezas: incluir más formulación o más gráficos, ser más didáctico o exponer de un modo más formal y rápido, etc. Recuerda que solo tú sabes lo que te has estudiado y que contenido tiene tu tema, el Tribunal solo tiene el título y la orientación de contenidos de los secretarios del Tribunal, que han sido opositores. Esto quiere decir que los contenidos no están escritos en piedra, son flexibles.
    }
\vspace{1cm}

\modificaciones{
    
    }

\newpage
\MyTOC
\newpage

\mylistofequations
\clearpage

\MyListOfFigures
\clearpage


\pagenumbering{arabic}
\setcounter{page}{1}


\section{Introducción}



\subsection{Contextualización}


\subsubsection{Relevancia}




\subsubsection{Problemática}



\subsection{Estructura de la exposición}




























\clearpage
\section{Organización Multilateral del Comercio: Acuerdos Multiltaerales distintos a los de mercancias}
\puntosclave{
    El objetivo del GATS es liberalizar el comercio de servicios. y para ellos incluye una serie de compromisos generales y específicos que se aplican a los Miembros de la OMC.
    
    Abarca todos los sectores (con alguna excepción), es un acuerdo complejo y muy flexible pues los miembros deciden que liberalizar y en qué grado (AM y TN) en sus listas de compromisos específicos. y atendiendo a cada uno de los 4 modos de suministro. 
    
    GATS está sujeto a la cláusula de NMF. 
    
    Lo que restringe (y liberaliza) el comercio de servicios es la regulación, por eso liberalizar servicios es más complejo que eliminar aranceles a las mercancías. GA TS especifica que la regulación debe ser objetiva, trasparente, y proporcionada.  
    
    Los compromisos iniciales de GA TS no eran muy ambiciosos, pero evitan retrocesos proteccionistas. Los avances de la Ronda de Ooha también han sido escasos, con lagunas excepciones como el waiver de servicios. 
    
    Por la vía plurilateral se ha conseguido avanzar en reglamentación nacional.

    El objetivo del TRIPS establecer reglas comunes de protección de los DPI que son vinculantes para los Miembros de la OMC y que cuyas diferencias se pueden resolvermediante el OSO de la OMC. Se busca un equilibrio entre ofrecer incentivos a la innovación y no restringir el acceso a tecnología, especialmente en determinados casos como por ejemplo las patentes. por sus efectos sobre la salud pública 
    
    La protección se lleva a cabo mediante una serie de normas mínimas de protección respecto a siete tipos de DPI, y además se introduce normativa para evitar prácticas anticompetitivas, asegurar la aplicación y cumplimiento, y unos periodos de transición favorables para los PEO. 
    
    Las novedades más importantes de la Ronda de Doha respecto al TR1Ps se han producido en relación a las patentes, que han pennitido aprobar licencias obligatorias para la producción y exportación a países sin capacidad de producción de medicamentos, y más recientemente agilizar estas excepciones para el caso de vacunas COVID.
}


En primer lugar, vamos a revisar los principales acuerdos multilaterales relativos a materias distintas al comercio de mercancias. Estos acuerdos multilaterales son, por así decirlo, fruto de la actividad de las instituciones o foros multilaterales de negociación, es decir, la Organización Mundial del Comercio. Estos son esencialmente dos: el Acuerdo General sobre el Comercio de Servicios (GATS) y el Acuerdo relativo a los Derechos de Propiedad Intelectual relacionados con el Comercio (ADPIC).

El GATS constituye una pieza fundamental de la arquitectura comerial internacional. ¿Por qué? En primer lugar porque el sector servicios, a lo largo de las últimas décadas, ha superado con creces la importancia de cualquier otro sector productivo en las economías más avanzadas; representando cerca del $70\%$ de la producción y el empleo mundial.\footnote{Además, el comercio de servicios supone en torno al $20\%$ del comercio mundial y el $50\%$ del comercio en términos de valor añadido (TiVA). En España estos porcentajes son aún mayores por nuestra especialización relativa en servicios: casi un tercio del comercio bruto, y un $60\%$ en valor añadido. Debemos también tener en cuenta que es mucho más reiliente a eventos adversos, tanto en la política como en lo económico, que el comerci omundial de mercancias. 
} Pero también porque constituye el primer y único conjunto de reglas multilaterales que regulan el comercio de servicios. 

Por otro lado, el Acuerdo sobre los Derechos de Propiedad Intelectual relacionados con el Comercio (ADPIC; \textit{Trade Related aspects of Intellectual Property Rights} en inglés, TRIPS) es el resultado de un largo camino de reivindicaciones y evidencias que acabaron volteando la balanza a su favor en la Ronda de Uruguay. Desde una perspectiva local, es probable que el acuerdo TRIPS sea casi más importante que el Acuerdo GATS: según datos de la Comisión Europea los sectores de la economía intensivos en DPIs representan el 47\% del PIB de la UE-$27$, el $30\%$ del empleo y el $80\%$ del comercio (exportaciones e importaciones):\footnote{\href{https://www.oepm.es/es/detalle-noticia/Nuevo-informe-EUIPO-EPO-sobre-Industrias-intensivas-en-derechos-de-propiedad-intelectual-y-resultados-economicos-en-la-Union-Europea/}{\textit{Nuevo informe EUIPO-EPO sobre Industrias intensivas en derechos de propiedad intelectual y resultados económicos en la Unión Europea}. Oficina Española de Patentes y Marcas (2022)}
}

\begin{center}
\begin{tabular}{|p{0.40\linewidth}|p{0.30\linewidth}|p{0.20\linewidth}|}
\hline
\textbf{Sectores Intesivos en DPI} & \textbf{Valor Añadido/PIB}\par\textbf{(en millones de EUR)} & \textbf{Proporción del PIB total de la UE} \\
\hline
Intensivos en marcas & $5.217.903$ & $38,5\%$\\
\hline
Intesivos en dibujos y modelos & $2.101.305$ & $15,5\%$\\
\hline
Intensivos en patentes & $2.361.457$ & $17,4\%$ \\
\hline
Intensivos en derechos de autor & $934.176$ & $6,9\%$\\
\hline
Intensivos en IG & $15.011$ & $0,1\%$\\
\hline
Intensivos en DOV & $187.774$ & $1,4\%$\\
\hline
\textbf{Todos los sectores intensivos en DPI} & $6.375.796$ & $47,1\%$\\
\hline
PIB total UE & $13.541.581$ & $100\%$\\
\hline
\multicolumn{3}{c}{\scriptsize NOTA$1$.- Debido al solapamiento en el uso de DPI, la suma de las cifras de cada uno de los DPIs supera la cifra total correspondiente a los sectores intensivos en DPI}\\
\multicolumn{3}{c}{\scriptsize NOTA$2$.- Se trata de la contribución media al PIB UE-$27$ entre 2017 y 2019. Para actualizar los datos: Oficina Europea de Patentes y Oficina de Propiedad Intelectual de la UE}
\end{tabular}
\end{center}

Además, al igual que el GATS, es el primer Acuerdo en incoporar, por primera vez, normas relativas a la propiedad intelectual (PI) dentro de la arquitectura multilateral tradicional, lo que constituye un paso clave porque, si bien ya existían otros convenios, algunos de los cuales quedan incluidos en el ADPIC, al integrarse como acuerdo multilateral de la OMC se convierten en obligatorios para todos los miembros y están sujetos al Mecanismo de Solución de Diferencias de la OMC.\footnote{Como podemos observar, se ha evitado esgrimir las razones \textit{económicas} que justifican la protección de la propiedad intelectual puesto que es objeto de otros temas. Sin embargo, a modo de recordatorio, lo que suele esgrimirse principalmente es: (i) incentiva la innovación per se ya que la innovación deja de considerarse un bien público puro (no rival, excluible) y se convierte en un bien excluible; (ii) incentiva la innovación, ya que genera un monopolio temporal que otorga beneficios extraordinarios y permite recuperar los costes de la innovación; y, (iii) la innovación supone, a priori, elevados costes fijos y por tanto objeto de importantes economías de escala que, es posible, no podrían llegar a generarse en la estrechez del mercado doméstico.

Con todo, la justificación económica del ADPIC podría ayudar a comprender porque la balanza basculó a su favor en la Ronda de Uruguay. Sin embargo, la valoración de ambos tratados será objeto de crítica en el último apartado de la sección.
} 

\textbf{¿Qué queremos decir?} En primer lugar, por tanto, veamos las particularidades de cada tratado para poder comprender su alcance e importancia mejor. Posteriormente, una vez comprendamos sus características y objetivos, estaremos preparados parallevar a cabo una valoración conjunto de su desempeño, así como también de las críticas que se han arrojado sobre estos. 

\subsection{GATS: ¿Cuáles son sus objetivos?}
Como podemos intuir, el objetivo del GATS es similar al GATT: la \textbf{expansión del comercio para fomentar el crecimiento y el desarrollo económico}. Con esto en mente, sus dos principales funciones son:
\begin{la}
    \item Servir de foro de negociación para la liberalización del comercio de servicios, lo que se lleva a cabo en las rondas de negociación; y,
    \item Definir un conjunto de nonnas comunes para el comercio de servicios, lo que mejora la trasparencia y reduce la incertidumbre, ya que los agentes cuentan con un marco estable y predecible de nonnas que regulan el comercio de servicios a nivel multilateral. 
\end{la}

Sin embargo, existen dos diferencias esenciales con respecto al GATT. La primera, y más obvia, es su ámbito de aplicación: el GATS es un \textbf{tratado que versa sobre el comercio de servicios}, actividades complejas que no pueden restringirse mediante barreras normales (fronteras, aranceles, etc.). ¿Cuáes serán, por tanto, las barreras al comercio de servicios? Las barreras al comercio de servicios son fundamentalmente de carácter regulatorio (legales) y podemos distinguir dos categirías fundamentales:\footnote{No obstante, conviene no perder de vista otros elementos que dificultan la prestación de servicios a nivel internacional, aunque no constituyan barreras en sentido estricto, tales como: (i) costes de transacción y de información, asociados a diferencias legales, lingüísticas o culturales; (ii) calidad de las instituciones, derivada de la regulación; (iii) problemas de conectividad telefónica y digital; (iv) costes de transporte, si es necesario el desplazamiento; y, (v) costes arancelarios, por importación/exportación de mercancías necesarias.
}
\begin{lnum}
    \item \textbf{Acceso al Mercado}. Aquellas que restringen o limitan el acceso al mercado, como puede ser la prestación de servicios a través de una empresa pública estatal o la venta de medicamentos únicamente bajo licencia. 
    \item \textbf{Trato diferenciado}. Aquellas que implican un trato diferenciado entre el proveedor nacional y el extranjero. En este caso el proveedor extranjero puede acceder al mercado nacional pero en desigualdad de condiciones frente al proveedor local. Algunos ejemplos incluyen los requisitos de nacionalidad o de residencia. 
\end{lnum}

Además, existen otras barreras relacionadas con la concesión de licencias, como son la opacidad de requisitos para su obtención, incertidumbre en el tiempo de procesamiento de las solicitudes o la obligación de presentar la documentación en papel. Por tanto, liberar su comercio implica modificaciones del \textbf{ordenamiento jurídico interno} del país. Algo que, como se puede intuir, resulta especialmente sensible y complejo en especial cuando se trata de servicios tradicionales y fuertemente intervenidos, lo que acaba dificiultando y limitando las ambiciones liberalizadoras.

La segunda diferencia crucial es su \textbf{flexibilidad}. Pero, ¿sobre qué? Si recordamos, uno de los principios básicos de funcionamiento del GATT es el trato nacional. Esto no es solo una floritura: sin equiparación de trato, el comercio de mercancias sigue teniendo barreras comerciales, como es obvio. Ahora bien, en el GATS las cosas son más complejas:\footnote{El nivel de Acceso a Mercado también se negocia pero es un concepto exclusivo del GATS que no tiene sentido en el GATT
} el \textbf{Trato Nacional se negocia} e, incluso, se podrá renegociar.\footnote{
Evidetemente, esto debería traer consigo una debida compensación, bien sea en forma de compensación económica o de mayor liberalización en otro sector. Al fin y al cabo, al ser derecho dispositivo, no existe ninguna forma de \textit{renegociar} tasada, será la que mutuamente accedan a concederse las partes.
} en forma de mayor liberalización en otro sector. Es decir, los países pueden decidir: (i) si conceden o no Trato Nacional; (ii) en qué grado lo conceden; (iii) en qué sector lo hacen; y, (iv) en qué modo de suministro.  

La razón de esta flexibilidad es muy razonable. En primer lugar, evita interferir en los objetivos de política pública nacional. Como decíamos, algunos servicios, como la intermediación bancaria o la salud, son sectores altamente regulados que atienden a motivaciones, tradiciones y preferencias internas. De esta forma, es posible disponer de la flexibilidad suficiente como para alcanzar compromisos específicos sin perturbar el funcionamiento o pacto social interno. Por otro lado, esta flexibilidad busca garantizar que la normativa nacional no constituya un obstáculo al comercio. Es decir, al ser de carácter flexible, los miembros se ven impedidos de mantener regulaciones más discrecionales o poco transparentes, lo que acaba beneficiando indirectamente el comercio de servicios incluso en ausencia de liberalización específica.

La estructura del GATS es extremadamente sencilla. Consta de seis partes:
\begin{enumerate}[label=Parte \arabic*.-]
    \item \textbf{Alcance y Definición}. En esta parte se establece y determina el ámbito de aplicación del Acuerdo;
    \item \textbf{Obligaciones y disciplinas generales}. Se establecen las normas generales de actuación que obligan a todos los miembros, y son de aplicación en todo momento;
    \item \textbf{Compromisos específicos}. Se establecen la forma y modo a través del cual los socios contraen sus obligaciones específicas para con el resto de los miembros. Es decir, las normas que deben seguirse al liberalizar un determinado servicio;
    \item \textbf{Liberalización progresiva}. Se explica cómo realizar los compromisos, modificar las listas de compromisos, etc.;
\end{enumerate}

Las siguientes dos siendo una serie de disposiciones institucionales y finales, poco importantes para lo que nos concierne. Y, seguido, una serie de anexos sobre las exenciones aplicables a diferentes principios y las condiciones especiales relativas a determinados sectores. Veremos únicamente los puntos más importantes.

\subsubsection{Ámbito de aplicación (Parte I): ¿cuáles son los \textit{servicios}?}
En primer lugar, como ya hemos avanzado, el ámbito de aplicación. 

El GATS cubre todos los servicios y todas las medidas del sector público que afecten al comercio de servicios, salvo dos excepciones tasadas: (i) los Servicios Públicos, servicios suministrados por el Sector Público (directa o indirectamente) en ejercicios de la autoridad gubernamental;\footnote{Por clarificar, que un servicio sea proporcionado por el Sector Público en ejercicio de su autoridad implicaría que su provisión se da en condiciones: \textbf{no comerciales y sin competencia}. En el caso de España, econtraríamos dentro de este tipo los Servicios Públicos provistos únicamente por el Estado (muy pocos), por ejemplo, la Defensa Nacional y el Control Aereo. Pensemos, por ejemplo, que sanidad, educación o Seguridad Social tienen su homólogo privado, por lo que no quedarían estríctamente excluidos del GATS por esta cláusula.
} y, (ii) lo derechos de tráfico aéreo y servicios relacionados.

Cuando los miembros decidan que conceden trato nacional a una sector concreto, lo más importante sería determinar el modo en que debe ser suminitrado el servicio para disponer de esta ventaja. Esto responde a una razón muy sencilla: los servicios son, además de extremedamente heterogeneos, objetos comerciales mucho más complejos que las mercancias, que solo se pueden proveer por entrega/paso transfronterizo. Por tanto se distinguen, en este sentido, cuatro modos de provisión:
\begin{lnum}
    \item \textbf{Comercio transfronterizo}. Un servicio se presta mediante comercio transfronterizo siempre y cuando se preste a distancia entre dos estados miembros. El ejemplo paradigmático sería la provisión de servicios por vía electrónica o digital (consultoría o informes sobre investigación de mercado a través de internet), servicios postales o asesoramiento tele-médico; 
    \item \textbf{Consumo extranjero}. Un servicio se presta por consumo extranjero siempre y cuando sea el consumidor del servicio el que se desplaza a otro país para consumirlo. Como es evidente, el ejemplo clásico es el turismo. También serían los estudiantes en el extranjero, reparaciones de embarcaciones o aeronaves o los pacientes médicos que se trasladan al extranjero;
    \item \textbf{Presencia comercial}. Proveer el servicio mediante presencia comercial implica que el proveedor presta sus servicio en otro país en el que se establece para tal fin. Como es evidente, se trata del modo de suministro comunmente exigido (y prestado) pues está estrechamente ligado a la inversión extranjera directa y un conjunto de beneficios fiscales para el país.\footnote{De hecho, su importancia es tal que ha sido necesario desarrollar una estadística propia: \href{https://www.elibrary.imf.org/display/book/9781589061286/ch04.xml}{\textit{Foreign Affiliates Trade in Service} (FTAS)}, que recoge actividades de filiales extranjeras. La razón es que estos datos no quedan registrados en la Balanza de Pagos vía servicios, pues esta se guía por el criterio de residencia
    } Los ejemplos clásicos son la empresa multinacional, con sus filiales y/o sucursales;
    \item \textbf{Movimiento temporal de personas físicas}. Por último, estaría la provisión mediante desplazamiento de la persona trabajadora, no de la empresa. En estos casos, el traslado es temporal y no con la intención de residencia. Es un método muy importante en algunos sectores, como los servicios profesionales. Existe un interés específico por parte de los países en vías de desarrollo en este modo de prestación, ya que aún no teniendo capacidad de inversión, sí son competitivos en determinados servicios, generalmente de baja cualificación por su alta competitividad en costes. Algunos ejemplos serían los un gerente de hotel que se desplaza a uno de los hoteles del grupo en el extranjero, un instalador de un equipo de maquinaria especializada que se desplaza para instalar en el extranjero, o un profesor que da un curso en otro país. 
\end{lnum}

Podemos ver fácilmente los modos de prestación y las barreras que afectan a cada modo de prestación a partir de la siguiente infografía:

\imagenfit{Fig1.png}{Modos de prestación: Barreras principales}{\href{https://comercio.gob.es/comercio-internacional-servicios/Documents/251126_Barreras_acceso_comercio_servicios_2025.pdf}{\textit{Barreras de Acceso al Comercio de Servicios}. Secretaría de Estado de Comercio. Gobierno de España}}


\begin{specialblock}
    Los servicios que no quedan excluidos explícitamente del GATS son clasificados en doce sectores y más de cientocincuenta subsectores (\textit{Central Product Classification}, de Naciones Unidas). La clasificación propuesta, que no es obligatoria, sí es seguida por la mayoría de miembros. Se identifican las siguientes categorías: (i) Servicios a empresas (servicios profesionales, de consultoría); (ii) Servicios de comunicación (telecomunicaciones, servicios postales, etc.); (iii) Servicios de la construcción; (iv) Servicios de distribución; (v) Servicios de enseñanza; (vi) Servicios sociales y de salud; (vii) Servicios relacionados con el medio ambiente (viii) Servicios financieros; (ix) Servicios turísticos; (x) Servicios de esparcimiento, culturales y deportivos; (xi) Servicios de transporte; (xii) Otros servicios.

    \imagenfit{Fig2.png}{Categorías: Comercio de Servicios}{\href{https://comercio.gob.es/comercio-internacional-servicios/Documents/251126_Barreras_acceso_comercio_servicios_2025.pdf}{\textit{Barreras de Acceso al Comercio de Servicios}. Secretaría de Estado de Comercio. Gobierno de España}}
    
    Uno de los retos actuales a los que se enfrenta el GATS es la actualización del CPC. Aunque hay versiones posteriores, la última publicada por Naciones Unidad en 2015, se sigue utilizando mayoritariamente la versión provisional del 1991. El mayor problema es que el desarrollo de nuevas tecnologías y aparición de nuevos tipos de servicios hace que esta clasificación se haya quedado obsoleta.

    Como vemos, en realidad el GATS puede abarcar practicamente todos los servicios que pueden prestarse. Aunque debemos recalcar que, precisamente, la flexibilidad hace que entre los \textbf{servicios bajo el paraguas del GATS} y los \textbf{servicios efectivamente comprometidos} haya un verdadero abismo de diferencia, esto no ha obstado para que sea, precisamente, fuente de fuertes críticas que defienden la exclusión de determinados (o un gran número) de servicios. Quizás la más conocida es la mal denominada \href{https://www.realinstitutoelcano.org/analisis/la-excepcion-cultural/}{\textbf{excepción cultural}} según la cual, los sector ligados a la cultura deben ser considerados \textit{extra commercium} o, como mínimo, no meras mercancias. Esta excepción, tan alegada por Francia para el sector audiovisual (muy protegido en Francia y en general en la UE), es la razón por la que no hay compromisos en este sector por parte de la UE. Como sus demandas prosperaron, la actitud francesa ha sido la de rechazar cualquier tipo de concesiones en este ámbito y tratar de convencer a otros miembros para que hagan lo mismo; lo que hace que, obviamente, sea más dificil llegar a acuerdos.
\end{specialblock}

Sobre este ámbito amplio que acabamos de ver, los miembros de la OMC tienen una serie de compromisos generales y específicos.

\subsubsection{Obligaciones Generales: ¿cuáles son los compromisos?}
Ahora que conocemos su amplísimo, y tan ambiguo, ámbito de aplicación, ¿qué compromisos se adquieren en el marco del GATS?

Las obligaciones dimanantes del AGCS se pueden clasificar en dos grandes categorías: \textbf{obligaciones generales aplicables a todos los Miembros} y sectores de servicios, y \textbf{obligaciones que solo son aplicables a los sectores consignados en la lista de compromisos} de un Miembro. Dichos compromisos se recogen en las distintas listas, cuyo alcance puede variar ampliamente según los Miembros. Los términos y conceptos pertinentes son similares, aunque no necesariamente idénticos, a los utilizados en el GATT. Por ejemplo, el trato nacional constituye una obligación general en el comercio de mercancías en tanto que en el AGCS es una obligación negociable.\footnote{El AGCS permite, además, que en determinadas circunstancias los Miembros tomen o mantengan medidas que supongan un incumplimiento de sus obligaciones en el marco del Acuerdo, incluidos la prescripción de la nación más favorecida o los compromisos específicos. El artículo correspondiente abarca, entre otras, las medidas necesarias para: (i) proteger la moral pública o mantener el orden público; (ii) proteger la salud de las personas y de los animales, o preservar los vegetales; o, (iii) asegurar la conformidad con leyes y reglamentos que no sean incompatibles con el Acuerdo, tales como las medidas necesarias para impedir prácticas engañosas o fraudulentas.

Además, el Anexo sobre Servicios Financieros confiere a los Miembros el derecho de adoptar medidas por motivos cautelares, incluida la protección de inversores, depositantes, tenedores de pólizas o personas con las que un proveedor de servicios financieros tenga contraída una obligación fiduciaria, o para garantizar la integridad y estabilidad del sistema financiero. Por último, en caso de graves dificultades de balanza de pagos los Miembros pueden aplicar temporalmente una restricción del comercio, de manera no discriminatoria, a pesar de la existencia de compromisos específicos.
}

\paragraph{Trato de Nación Más Favorecida (art. 11): Multilateralidad}
El art. $2$ del AGCS dispone que cada Miembro debe otorgar inmediata e incondicionalmente a los servicios y a los proveedores de servicios de cualquier otro Miembro \textbf{un trato no menos favorable que el que conceda a los servicios similares y a los proveedores de servicios similares de cualquier otro país}. Esta disposición equivale a una prohibición, en principio, de acuerdos preferenciales entre grupos de Miembros en determinados sectores o de disposiciones de reciprocidad que limitan los beneficios de la adhesión a los socios comerciales que otorgan un trato similar.

Por un lado, esto implica que cualquier ventaja comercial concedida en el comercio de servicios con un país se debe multilateralizar (extender) al resto. Pero también hay que tener en cuenta que genera problemas de \textit{free rider}, puesto que al no haber compromisos ex ante, los países que no realizan concesiones se benefician igual de las que realizan el resto, lo que afecta al nivel de ambición de las negociaciones y, como es razonable, conduce a que haya más reticencias a conceder concesiones.

No obstante, existen también excepciones en forma de las denominadas exenciones de las obligaciones del artículo $2$. Se permitió que los Miembros solicitaran esas exenciones antes de que entrara en vigor el Acuerdo. Solo pueden concederse nuevas exenciones a los nuevos Miembros, en el momento de su adhesión, o, en el caso de los Miembros actuales, mediante la concesión de una exención. Todas las exenciones están sujetas a examen y en principio no deben exceder de un plazo de $10$ años. Además, el AGCS permite a grupos de Miembros la adhesión a acuerdos de integración económica o el reconocimiento mutuo de normas de reglamentación, certificados y similares si reúnen ciertas condiciones.

La más importante, como es normal, aplica sobre a acuerdos de integración económica (art. 5): \textbf{los acuerdos de integración económica están excluidas de la obligación de extender las ventajas del acuerdo a otros países que no formen parte del acuerdo regional}. No obstante, para que sea aplicable dicha excepción, los acuerdos deben cumplir dos condiciones: (i) carácter liberalizador, fomentando el comercio entre los miembros \textbf{sin aumentar restricciones frente a terceros};\footnote{La justificación teórica de esta condición proviene de las Teorías de KEMP y WAN. Proponen que el, para que un AEC sea  óptimo de PARETO, debe crear comercio dentro de la Unión Aduanera y evitar la desviación de comercio al resto del mundo (no es más restrictivo frente a países terceros).
} y, (ii) de cobertura amplia, que cubra un número suficiente de sectores (no sectoriales), modos de suministro y volumen de comercio.  de cobertura amplia (no sectoriales).\footnote{La justificación teórica de esta condición se encuentra en las teorías de Unión Aduanera de los hermanos WONNACCOT y las aportaciones de VINER, que consideran que siempre y cuando el acuerdo sea lo suficientemente amplio, es posible que los efectos positivos sobre las exportaciones (efecto creación de comercio) compensen los negativos sobre las importaciones (efecto desviación de comercio).

En la práctica se aplica una regla no estricta que entiende cumplido este requisito cuando se cubre el $90\%$ del comercio bilateral
}

\paragraph{Transparencia (art. 5): Ley y jurisdicción}
Por otro lado, los Miembros se obligan activamente a publicitar sus decisiones internas. La primera, y más importante, obligación que se desprende del artículo es la publicación con prontitud, salvo emergencia, de todas las medidas pertinentes de aplicación general que se refieran al presente Acuerdo o afecten a su funcionamiento. Así como también los acuerdos internacionales que se refieran o afecten al comercio de servicios y de los que sea signatario un Miembro. 

Por otro lado, siempre que se lleve a cabo una modificación que afecte significativamente al comercio de servicios abarcado por sus compromisos específicos en virtud del presente Acuerdo, o de la introducción de modificaciones en los ya existentes, se informará con prontitud, y por lo menos anualmente, al Consejo del Comercio de Servicios de su establecimiento (leyes, reglamentos o directrices administrativas). Paralelalemente, cada Miembro debe responder a las peticiones de información específica formuladas por los demás Miembros acerca de cualesquiera de sus medidas de aplicación general o acuerdos internacionales, y establecerá uno o más servicios encargados de facilitar información específica a los otros Miembros que lo soliciten sobre todas esas cuestiones. 

Otras obligaciones de aplicación general son el establecimiento de procedimientos de revisión administrativa y de recurso y de disciplinas aplicables a las actividades de monopolios y proveedores exclusivos.

Como se puede ver, el GATS no es un acuerdo ni completo ni cerrado, quedan muchas cosas por hacer. En particular, mejorar sus desarrollos en ámbitos como la reglamentación nacional, las compras públicas, las medidas de salvaguardia y las subvenciones.\footnote{Estos temas se han venido negociando en la OMC sin que hasta la fecha haya habido progresos sustanciales en el plano multilateral, pero sí se han producido avancen en algunas de ellos m ediante las in iciativas Conjuntas Multilaterales (por ejemplo, reglamentación nacional).
}

\subsubsection{Obligaciones específicas: ¿Que es la \textit{lista positiva}?}
Conocidas las obligaciones generales adquiridas por el simple hecho de pertenecer al GATS, veamos las obligaciones específicas que pueden derivarse del tratado. 

Todo miembro de la Organización \textbf{debe presentar una lista de compromisos específicos de liberalización}, aunque son libres de decidir en qué sectores y en qué grado de liberalización, lo que da flexibilidad al acuerdo. Estos compromisos se enumeran en la lista de compromisos de cada país, que hoy en día es una lista positiva: \textbf{solo está liberalizado lo que se incluye en la lista}. 

\begin{specialblock}{Sistema de \textbf{lista positiva}: ¿Cómo funciona?}
    Existen otros formatos de lista, principalmente tres alternativas:
    \begin{lnum}
        \item \textbf{Positiva} (GATS): solo está liberalizado lo que aparece en la lista. Primero se listan los sectores en los que se quiere hacer compromisos, y luego se explican los compromisos;
        \item \textbf{Negativa} (USMCA): se lista lo que no está liberalizado. Se incluyen solo los sectores en lo que no se realizan compromisos, es decir se listan solo las barreras al Trato Nacional y al Acceso a Mercado, y lo que no está consolidado;
        \item \textbf{Mixta} (TISA): lista positiva para Acceso a Mercado, negativa para Trato Nacional.
    \end{lnum}

    En el caso del GATS, con lista positiva, la liberalización (o su grado) se negocia o compromete por \textbf{sectores}. Es decir, en su lista, el país refleja/incluye una serie de compromisos horizontales que se aplican a todos los sectores que se incluyan en la lista, y unos compromisos específicos para determinados sectores que se incluyen en la lista y que solo se aplican a dicho sector. Además, para absolutamente todos los sectores incluidos en la lista, el país debe especificar el grado de liberalización. En este sentido, existen dos grandes antagonistas:
    \begin{la}
        \item \textbf{Liberalización nula}. Si el país no pretende liberalizar el sector, pondría [S/C: \textit{Sin consolidar}]. Esto significa que el país, para ese modo y sector, se reserva el derecho de tomar cualquier medida que vaya en contra de: (i) el Acceso a Mercado, si la señal está en esa columna; o, (ii) el Trato Nacional, si el S/C está en esa columna. Es decir, el país se reserva el derecho de modificar su reglamentación y puede cambiar las condiciones de entrada a los proveedores; o,
        \item \textbf{Liberalización total}. Si el país especifica [\textit{Ninguna}] en la lista, esto implica no se reserva el derecho a imponer ninguna medida restrictiva ni de Acceso a Mercado ni de Trato Nacional. 
    \end{la}
    
    Existe la posibilidad de observar un término medio: se puede especificar cualquier tipo de limitación concreta al Acceso a Mercado o al Trato Nacional, siempre y cuando se incluya en la lista. Por ejemplo, \textit{máximo 100 licencias de banca minorista}, o \textit{máximo participación extranjera $49\%$}. Este es un sencillo ejemplo de esto, que nos permite ver como sería una lista positiva en el marco del GATS. (Podemos ponerlo en la pizarra). 

    \begin{center}
    \begin{tabular}{|p{0.25\linewidth}|p{0.25\linewidth}|p{0.25\linewidth}|p{0.1\linewidth}|}
    \hline
    \textbf{Sector o subsector} & \textbf{Limitaciones al Acceso a Mercado} & \textbf{Limitaciones al Trato Nacional} & {\scriptsize Compromisos adicionales}\\
    \hline
    \textbf{Compromisos horizontales} & 1) Sin Consolidar \par 2) Ninguna\par 3) Sin consolidar\par 4) Sin consolidar & 1) Sin Consolidar \par 2) Ninguna \par 3) Sin consolidar \par 4) Sin consolidar & \\
    \hline
    \textbf{Compromisos específicos}\par\par\par
    Banca Comercial\par\par Servicios Financieros & 1) Ninguna\par 2) Ninguna\par 3) Participación extranjera máxima del $49\%$ en el capital social\par 4) Sin consolidar & 1) Ninguna\par 2) Ninguna\par 3) Participación extranjera máxima del $49\%$ en el capital social\par 4) Sin consolidar & \\
    \hline
    \end{tabular}
    \end{center}

    Como vemos, tendrá cuatro columnas, en las que se especifica: el sector y subsector, las limitaciones al Acceso a Mercado, las limitaciones al Trato Nacional y una última columna de compromisos adicionales. En las filas, por otra parte, se incluyen los compromisos horizontales y los específicos de cada sector. Los números del $1)$ al $4)$ hacen referencia a los compromisos en cada \textbf{modo de suministro}.\footnote{Esto es, en cualquier caso, imperativo: se debe aclarar el tratamiento que se va a dar en los cuatro modos de provisión.
    } En este ejemplo, los compromisos adicionales liberalizan la prestación de servicios en modo dos y en los compromisos especificos del subsector de banca comercial, dentro del sector de servicios financieros, liberalizan adicionalmente el modo uno. Además, se especifica una limitación al capital extranjero, que es una barrera de Trato Nacional, y será la única que se permitirá en el subsector porque aparece en la lista proporcionada.

    Como es evidente, las listas varían de país a país atendiendo a sus preferencias (compromisos específicos). La complejidad de estas listas puede también ser radicalmente diferente. Por ejemplo, si un país decide abarcar compromisos horizontales en muchos sectores ($12$) y un buen número de subsectores, puede llegar a tener cientos de páginas. Sin embargo, si el país decide liberalizar un sector concreto completa y únicamente, puede ser perfectamente de un único folio.
\end{specialblock}

Estos compromisos y restriccione que se adquieren y negocian voluntariamente son de dos tipos.

\paragraph{Acceso a Mercado}
En primer lugar, los miembros pueden comprometerse a no aplicar restricciones a la provisión de servicios vía comercio de servicios. Es decir: \textbf{permitir el acceso a su mercado}. 

Como existe un amplio abanico de restricciones posible, el GATS especifíca esencialmente seis:
\begin{lnum}
    \item Limitación al número de proveedores de servicios. Por ejemplo, los bancos o las compañías de seguro extranjeros sólo podrán establecer un número determinado de filiales o sucursales;
    \item Limitación del valor de los activos o transacciones de servicios. Por ejemplo, sólo podrá colocarse en compañías extranjeras el $10\%$ del valor del reaseguro;
    \item Limitación del número de operaciones de servicios o de la cuantía total de la producción de servicios;
    \item Limitación del número de personas fisicas que puedan emplearse en un determinado sector de servicios o que un proveedor de servicios pueda emplear. Por ejemplo, la mayoría de los miembros del consejo de administración deben ser ciudadanos del país; y,
    \item Limitación de la participación de capital extranjero. Por ejemplo, una participación máxima del $49\%$.
\end{lnum}

Por último, estarían las medidas que restrinjan o prescriban los tipos específicos de persona jurídica por medio de los cuales pueda prestarse un servicio. Por ejemplo, en la banca o los seguros, las filiales deben estar constituidas en sociedades). 

\paragraph{Trato Nacional (art. $3$)}
Por otro lado, los miembros pueden comprometerse a \textbf{dar el mismo trato a los proveedores nacionales y extranjeros}. El GATS lo define como aquel trato que asegura  que las condiciones de competencia no sean alteradas en favor del proveedor de servicios nacional. 

Se trata de medidas discriminatorias que no se recogen de forma explícita en la lista, pero que pueden darse en diferentes formas, como por ejemplo, requisitos de residencia para ejercer una profesión. La obligación de trato nacional no significa que las autoridades tengan que otorgar un trato idéntico a proveedores de servicios nacionales y extranjeros, lo que importa es que les ofrezcan las mismas posibilidades de competir.

\subsection{TRIPS: ¿Cuáles son sus objetivos?}
Ahora que conocemos qué es, qué obetivos persigue y cómo funciona el GATS, podemos presentar el segundo Acuerdo Multilateral no relativo a las mercancias: el ADPIC (Acuerdo relativo a los Derechos de Propiedad Intelectual relacionados con el Comercio). En primer lugar, debemos saber que la propiedad intelectual se define como el \textbf{conjunto de derechos que corresponden a los autores y a otros titulares} por las obras y otros trabajos que sean fruto de su creación. Con esta definición en mente, el Acuerdo se marcó perseguir los siguientes objetivos:
\begin{la}
    \item \textbf{Marco normativo mínimo común: MCM}. Establecer unas normas mínimas de protección de los derechos de propiedad intelectual, que generen mayor certidumbre a los agentes que participan en el comercio internacional, y luchar contra las falsificaciones.\footnote{Cada año se venden en el mundo productos falsificados por valor de 461.000 Millones €, lo que supone el 2,5\% del comercio mundial.
    } y,
    \item \textbf{Soluciones conflictos}. Poner a disposición un mecanismo para la solución de diferencias, el de la OMC. 
\end{la}

\subsubsection{Antecedentes: WIPO}
A lo largo de la segunda mitad del s.XX, las ideas y los conocimientos empezaron a adquirir un papel cada vez más relevante en el comercio. La mayor parte del valor de los medicamentos y otros productos de alta tecnología empezaba a atribuirse cada vez más a la cantidad de invención, innovación, investigación, diseño y pruebas que requerían. En la misma linea, películas, grabaciones musicales y libros, y posteriormente programas de ordenador y servicios en línea se combranan y vendían (al igual que hoy) por la información y creatividad que contienen. Por último, muchos productos que solían ser objeto de comercio como productos de baja tecnología empezaron a contener una porción mayor de su valor atribuible al diseño y marca que ostentaban, por ejemplo, las prendas de vestir de marca o las obtenciones vegetales.

Los creadores tenían el derecho de impedir que otros utilizaran sus invenciones, diseños o demás creaciones: los \textbf{derechos de propiedad intelectual}.\footnote{Los derechos depropiedad intelectual revisten diversas formas: libros, pinturas y películas quedan protegidos por el derecho de autor; las invenciones pueden patentarse; los nombres comerciales y los logotipos de productos pueden registrarse como marcas de fábrica o de comercio; y así sucesivamente. Los gobiernos y los parlamentos han conferido a los creadores esos derechos como incentivo para generar ideas que beneficien a la sociedad en su conjunto.
} Sin embargo, el grado de protección y observancia de esos derechos variaba considerablemente en los distintos países del mundo y, a medida que la propiedad intelectual iba adquiriendo mayor importancia, esas diferencias se convirtieron en una fuente de tensiones en las relaciones económicas internacionales. (Principalmente porque los acuerdos se circunscribían a los países desarrollados).

El primer intento por hacer valer esta protección de manera mucho más amplia se obtuvo con la creación de la WIPO (World Intelectual Property Organization). En el marco de esta organización, se desarrollaron una serie de convenciones específicas que regulaban los derechos de propiedad intelectual a nivel multilateral (entre miembros).\footnote{Fundamentalmente, la Convención de París (referida a la protección de patentes, dibojos y modelos industriales, etc.) y la Convención de Berna para la protección de obras artísticas y literarias (es decir, derechos de autor).
} Sin embargo, a finales de los años 70's, los países desarrollados (y EEUU en particular) vieron el sistema infusificente para proteger adecuadamente los intereses de sus industrias intensivas en conocimiento. Y, si bien el deseo de protección era tal vez excesivo, ciertamente la WIPO tenía poca capacidad para hacer cumplir incluso los acuerdos y para actuar como mecanismo de solución de diferencias entre sus miembros. Al mismo tiempo, en los años a finales de los años 70's y principios de los años 80's, la piratería comienza a crecer como resultado del intento de algunos países en desarrollo de intensificar su proceso de industrialización y acortar distancias respecto a los países desarrollados mediante la transferencia forzoso de conocimiento y la ingenieria inversa. 

Durante la Rondade Tokio (1973), se busca alcanzar un acuerdo en materia de Derechos de Propiedad Intelectual, a instancia de la colaboración del Gobierno estadounidense y un grupo selecto de multinacionales, con la creación del \textit{Anticounterfeiting Coalition} (coalición para la negociación), cuyo resultado fue la propuesta de Código de falsificación. Aunque la inciiativa no prosperó, en 1979 los Estados Unidos y la Comunidad Europea alcanzan un acuerdo para presentar conjuntamente una propuesta de Acuerdo sobre Medidas para Desincentivar la Importación de Bienes Falsificados, que tampoco vería su aprobación con el final de la Ronda de Tokio (1979). 

\textcolor{red}{Moving (or threatening to relocate) the matter to another body where strong countries are more likely to get their way is one method of exiting a situation. For instance, the European Community and the United States moved the issue to the GATT in the early 1980s when they were unable to secure the necessary majority in the World Intellectual Property Organization for broader intellectual property protection, and they were able to reach an agreement on the Trade-Related Aspects of Intellectual Property Rights (TRIPS) Agreement in 1994.}

En vistas de la siguiente Ronda de negociaciones, Estados Unidos, apoyado por Japón, presentó una nueva propuesta en el Comité Preparatorio, en 1986. El objetivo principal sería tratar los derechos de propiedad intelectual en el ámbito del GATT en general, y no sólo los relativos a marcas falsificadas. Con su influencia, a pesar de las reticencias de los países en desarrollo, la Declaración Ministerial de 1986 materializó las aspiraciones estadounidenses (y de, en general, el bloque más avanzado) y acordaron incluir el programa en las próximas negociaciones del GATT.\footnote{Este declaración acordaba un marco de negociación que incluyera tanto normas sustantivas relativas a los Derechos de Propiedad Intelectual, como procedimientos particulares de aplicación y solución de diferencias. 
} A partir de entonces, se inicia la trayectoria que dará lugar al ADPIC. En particular, se da inicio a la redacción de los textos, prestando consideración a las propuestas de los miembros, ya que las posiciones eran muy divergentes. 

En diciembre de 1991, el Director General del GATT (DUNKEL) presentó un texto de compromiso del Proyecto de ADPIC. Después de esta presentación, sufriría pequeñas modificaciones hasta su adopción final en Marrakech, en abril de 1994. 

Por tanto, la Ronda Uruguay fue la que lo consiguió. El Acuerdo de la OMC sobre los ADPIC constituye un intento de reducir las diferencias en la manera de proteger esos derechos en los distintos países del mundo y de someterlos a normas internacionales comunes. En él se establecen niveles mínimos de protección que cada gobierno ha de otorgar a la propiedad intelectual de los demás Miembros de la OMC. Al hacerlo, establece un equilibrio entre los beneficios a largo plazo y los posibles costos a corto plazo resultantes para la sociedad. Los beneficios a largo plazo para la sociedad se producen cuando la protección de la propiedad intelectual fomenta la creación y la invención, especialmente cuando expira el período de protección y las creaciones e invenciones pasan a ser del dominio público. Los gobiernos están autorizados a reducir los costos a corto plazo que puedan producirse mediante diversas excepciones, por ejemplo hacer frente a los problemas relativos a la salud pública. Y actualmente, cuando surgen diferencias comerciales con respecto a derechos de propiedad intelectual, puede recurrirse al sistema de solución de diferencias de la OMC. 

\paragraph{Ronda de Uruguay: ADPIC}
El Acuerdo de la OMC sobre los Aspectos de los Derechos de Propiedad Intelectual relacionados con el Comercio (Acuerdo sobre los ADPIC) es el acuerdo multilateral más completo sobre propiedad intelectual, fundamental para facilitar el comercio de conocimientos y contenidos creativos, solucionar diferencias comerciales relacionadas con la propiedad intelectual y dar a los Miembros de la OMC margen para lograr sus objetivos de política nacionales. 

Para alcanzar estos objetivos, el Acuerdo TRIPS aprovecha los principales convenios y convenciones sobre los derechos de propiedad intelectual incorporando (por remisión) la mayoría de las disposiciones de éstos. Y, por otro lado, prevé que los países podrán, en cumplimiento de estos convenios, garantizar una protección más amplia de lo exigida por el propio Acuerdo, siempre que tal protección no infrinja las disposiciones del mismo. Abarca cinco amplias cuestiones que dan lugar, al mismo tiempo, a su propia estructura:
\begin{lnum}
    \item \textbf{Principios y obligaciones generales}. Cómo deben aplicarse los principios básicos del sistema de comercio y otros acuerdos internacionales sobre propiedad intelectual;
    \item \textbf{Normas mínimas de protección} (establece siete categorías). Cómo prestar protección adecuada a los derechos de propiedad intelectual;
    \item \textbf{Prácticas anticompetitivas en las licencias contractuales}. Cómo deben los países hacer respetar adecuadamente esos derechos en sus territorios;
    \item \textbf{Medidas nacionales para asegurar la aplicación y cumplimiento de los DPI}. Cómo resolver las diferencias en materia de propiedad intelectual entre Miembros de la OMC; y,
    \item \textbf{Periodos de carencia de las reglas en ámbito nacional}. Disposiciones transitorias especiales durante el período de establecimiento del nuevo sistema.
\end{lnum}

\subsubsection{Principios básicos: ¿Qué se considera normal en el ámbito IPO?}
Al igual que en el GATT y en el AGCS, el punto de partida del Acuerdo sobre la propiedad intelectual son los principios básicos.

\paragraph{Principio de no discriminación: MFN y TN}
Reviste especial importancia el \textbf{principio de no discriminación}.

Como sabemos, este está compuesto de dos obligaciones: (i) \textbf{trato nacional}, igualdad de trato para nacionales y extranjeros; y, (ii) trato de la \textbf{nación más favorecida}, igualdad de trato para los nacionales de todos los interlocutores comerciales en el marco de la OMC. 

\paragraph{Protección Equilibrada}
Además, existe un principio adicional y particularmente importante en este acuerdo: el p\textbf{rincipio de protección equilibrada}.

El ADPIC establece que la propedad intelectual debe contribuir a la innovación y a la transferencia de tecnología, deben beneficiarse tanto produtores como usuarios, y debe contribuir a un mayor bienestar económico y social. La clave aquí está en buscar el equilibrio entre: (i) los beneficios a largo plazo, el incentivo a la innovación que proporciona el monopolio temporal y la transferencia de tecnología cuando la patente o el DPI que sea expira; y, (ii) los costes a corto plazo, porque dificulta la transferencia de tecnología ya que excluye del acceso a la tecnología, y aumenta los costes para todo el que no es propietario del derecho. 

Un ejemplo claro, y también fuente de numerosas controversias es el ámbito de la industria farmaceútica: los países con menos recursos no se pueden ver impedidos de acceder a medicamentos especialmente importante para proteger su desarrollo y capital humano al ser los más afectados por enfermedades especialmente severa, como el VIH.

Esto obliga a buscar un equilibro en: (i) el periodo de protección, que debe ser suficientemente largo para incentivar la innovación, pero no demasiado para que no evite la transferencia de tecnología; y, (ii) el grado de protección, especialmente en los países en desarrollo y los menos avanzados.

\subsubsection{Obligaciones específicas: ¿Cómo alcanzar esta protección?} 
Dicho equilibro se alcanza estableciendo un balance entre normas mínimas de protección y el expreso control de posibles prácticas anticompetitivas que puedan ocasionarse. 

Empecemos por esta segunda.

\paragraph{Prácticas anticompetitivas}
Las disposiciones en contra de prácticas anticompetitivas persiguen impedir que los titulares del derecho abusen de el y acaben restringiendo la competencia o la transferencia de tecnologías. Por ejemplo, otorgando licencias en términos abusivos.

\paragraph{Normas mínimas de protección}
En segundo lugar, como hemos dicho, tendríamos la protección efectiva de los derechos. El ADPIC establece unos estándares mínimos universales para un total de siete tipos de de derechos: un país pude proteger más en su legislación, pero no puede proteger menos. Brevemente, cada uno de los sectores contempla las siguiente normas mínimas:
\begin{la}
    \item \textbf{Patentes}. Las patentes se otorgan a las invenciones por un mínimo de $20$ años. Para ser considerada como invención debe cumplir dos requisitos: \textbf{ser nueva y tener aplicación industrial}. El Acuerdo dispone que se otorgarán patentes por las invenciones en todos los campos de la tecnología, tanto de productos como de procesos, pero quedan exluidos: las invenciones contrarias al orden público, moral, salud o vida de los seres vivos, o el medio ambiente; los métodos de diagnóstico, terapéuticos y quirúrgicos para el tratamiento de personas o animales; las plantas y los animales, pero sí los microorganismos. La petente confiere la propiedad exclusiva para la \textbf{producción, venta y distribución}, de forma que obliga a quien quiera producir/importar/vender esas invenciones a solicitar una licencia o autorización del titular de la patente.\footnote{Existen algunas excepciones a las licencias obligatorias de algunas patentes para evitar el abuso. En estos casos, lo habitual es que se obligue la concesión de una licencia cuando el fabricante haya intentado, sin éxito, obtener la autorización correspondiente (licencia) en términos y condiciones razonables (esto es una aplicación de la obligación general de evitar prácticas anticompetitivas o abusivas que hemos estudiado antes). El caso más emblemáticos es el de los medicamentos.
    } Sin embargo, también impone la obligación de divulgación: los titulares de patentes están obligados a divulgar la información técnica sobre esas invenciones, para que esta pueda ser aprovechada para nuevas investigaciones o para poder aplicar la invención una vez expirada la patente;
    \item \textbf{Derechos de autor}. El derecho de autor protege las obras en el campo de la literatura, la ciencia, el arte u otra forma de expresión, siempre que esa obra sea una \textbf{creación original del autor}, por un mínimo de $50$ años. Confiere derechos de reproducción, derechos de interpretación, derechos de grabación, derechos de cinematografia, derechos de radiodifusión y derechos de traducción y adaptación. Además, las leyes confieren a los autores derechos morales que les permite reivindicar la autoría de la obra y oponerse a toda distorsión que pudiera redundar en menoscabo de su reputación u honor;
    \item \textbf{Marcas}. Una marca es un signo que sirve para distinguir los productos o servicios de una empresa industrial o comercial de los de otras empresas. Tiene doble finalidad. Por un lado, busca ayudar a sus propietarios a vender y promocionar sus productos y mantener o mejorar la calidad de los mismos y, por otro lado ayudar al consumidor a elegir entre varias posibilidades o identificar un producto. La duración del primer registro es como mínimo de siete años pero puede renovarse indefinidamente. Su registro da derecho exclusivo de uso, si bien dejará de estar protegida (y podrá anularse) después de un periodo ininterrumpido mínimo de tres años en desuso. Se pueden conceder licencias para el uso de marcas, pero a diferencia de las patentes, no se permiten las licencias obligatorias de marcas. El titular de la marca tiene derecho a cederla con o sin la transferencia de la empresa a que pertenezca la marca;
    \item \textbf{Diseños industriales}. Los dibujos y modelos industriales deben ser nuevos u originales. Se protegen sobre todo los productos de consumo, de los que los textiles, el cuero y los productos de cuero y los automóviles son los ejemplos más corrientes. La duración mínima de protección es de 10 años. El titular tiene derecho exclusivo de uso y puede impedir que terceros, sin su consentimiento, fabriquen, vendan o importen artículos que ostenten o incorporen un dibujo o modelo que sea una copia del dibujo o modelo protegido;
    \item \textbf{Indicaciones geográficas}. Las indicaciones geográficas (IGs) identifican un producto con un nivel de calidad, reputación u otra característica imputable a su origen geográfico. En los países desarrollados estas leyes se aplican predominantemente a los vinos y aguardientes (champagne, Coñac, Jerez) pero existen otros productos (Rochefort, Basmati). El Acuerdo establece que se deben aplicar medidas para impedir el uso de IGs falsas que induzcan a error por parte de los consumidores, o cuyo objetivo sea la competencia desleal.\footnote{El ejemplo más conocido es el del champagne, término asociado al vino producido en una determinada región de Francia: no es lícito llamar champagne al vino producido en otras partes (Argentina o Estados Unidos) aunque en el país productor ese vino se considere comparable al champagne francés. De ahí que el homólogo español se llame \textit{cava}. En España suele conocerse este derecho como \textbf{denominación de origen}. Para el caso de licores, la protección es incluso mayor: en la Ronda Uruguay se acordó un mandato de negociación para estos que otorgara mayor protección y ayudara a evitar la usurpación de sus nombres. En estos momentos la negociación continúa. (Como anécdota, la Ginebra es de los pocos licores que no tiene denominación de origen y su distribución bajo ese nombre no está ligado a ninguna demarcación geográfica).
    }
    \item \textbf{Trazado de circuitos integrados}. La base de la protección prevista para los esquemas de trazado (\textbf{topografías}) de los circuitos integrados es el Tratado de Washington sobre la Propiedad Intelectual respecto de los Circuitos Integrados, concluido en el marco de la Organización Mundial de la Propiedad Intelectual. Se adoptó en 1989, pero aún no ha entrado en vigor. En el Acuerdo sobre los ADPIC se añaden una serie de disposiciones: por ejemplo, la protección debe otorgarse por un plazo mínimo de $10$ años.
    \item \textbf{Información no divulgada} (secretos comerciales). Los secretos comerciales y otros tipos de información no divulgada que tengan valor comercial deben estar protegidos contra abusos de confianza y otros actos contrarios a los usos comerciales honestos. Ahora bien, deben haberse adoptado medidas razonables para mantener secreta la información. También deben estar protegidos contra todo uso comercial desleal los datos de pruebas facilitados a los gobiernos con el fin de obtener autorización para la comercialización de productos farmacéuticos o productos químicos agrícolas nuevos. En la práctica se suele usar cuando no se quiere revelar la información (como ocurre con una patente, el ejemplo clásico es Coca-Cola) o cuando algo no es patentable.
\end{la}

Como vemos, para cada una de las categorías se especifica la materia protegida (definición) y la duración mínima (salvo para marcas, la información no divulgada y las indicaciones geográficas, que tienen duración ilimitada), los derechos conferidos y las excepciones a esos derechos. 

\subsection{Valoración: ¿?}
El GATS
2. Para los países desarrollados (PO), la liberalización de los servicios es un claro interés ofensivo (es decir, interés por liberalizar mercados exteriores), ya que son economías muy terciarizadas y tienen ventaja comparativa en la provisión de muchos servicios. 
3. Los países en desarrollo (PED), por el contrario, tienen en general intereses defensivos. La liberalización de los servicios es al mismo tiempo un reto y una oportunidad. 
o Reto, porque implica hacer frente a una fuerte competencia internacional. 
o Oportunidad porque: • hacer concesiones en liberalización de servicios facilita las negociaciones agrícolas. en las que los PEO tienen mayor interés. determinadas formas de prestación, como el modo 4 de prestación de servicios mediante desplazamiento temporal de personas fisicas, sí son de interés para los PED. 
4. También se puede argumentar (valido para PD y PED) que la liberalización de servicios pennite acceso a servicios eficientes que son inputs esenciales para otras actividades económicas, y por tanto mejoran la productividad y competitividad: telecomunicaciones, servicios financieros, servicios TIC.

El TRJPS es posiblemente el acuerdo de la OMC que mayor controversia ha generado, por 2 razones.

1 1 1 . Es un claro interés ofensivo de los países desarrollados, que son los principales productores, y por tanto son sus empresas las que dejan de percibir rentas/royalties, además de perder competitividad como consecuencia de la falta de respeto de dichos derechos (recordemos que EEUU ha sido el principal impulsor del acuerdo). Por otro lado, perjudica a los países en desarrollo y PMAs, que no tienen ventajas comparativas en el desarrollo de este tipo de intangibles, y se ven obligados a pagar un precio mayor cuando los derechos están protegidos.

2. Los cortos períodos de transición permiten afianzar rápidamente los beneficios a los países desarrollados, porque el resto (PED. PMAs) no tienen tiempo a desarrollar la tecnología (o copiarla). a. En el caso de los medicamentos, como veremos, se han introducido mayores flexibilidades para paliar este problema, ya que en este caso se prima la salud ybienestar en los PED y PMAs: se permite la producción y distribución de genéricos de medicamentos en PED, lo que se consigue con concesión de licencias obligatorias.

medidas nacionales para asegurar la aplicación y cumplimiento de los DPI. Casos como el de China son un claro ejemplo de que la existencia de normas es condición necesaria pero no suficiente para el respeto de los DPI: casi tan importante como que existan buenas leyes es que éstas se hagan cumplir. Es necesario que los Gobiernos se aseguren: • (ex - ante) que un titular puede hacer valer sus derechos ante un tribunal • (ex- post) que las sanciones con suficientemente elevadas para desincentivar al infractor. Para ello, el TRIPS establece que: Acciones civiles: • Los procedimientos deben ser sencillos, justos y rápidos, y a un coste razonable. • Se podrán tomar medidas provisionales rápidas y eficaces para preservar las pruebas de una infracción, y para evitar que se produzca una infracción. • De hecho, muchas de las medidas están enfocadas a evitar que las mercancías entren en los circuitos comerciales, por ejemplo evitar la importación en aduana de mercancías sospechosas. • Una vez comprobada la infracción de un DPI, los tribunales pueden ordenar al infractor que pague al titular del derecho un resarcimiento para compensar el daño.  Acciones penales • Los tribunale�� nacionales deben estar capacitados para: o Destruir mercancía falsificada o piratería o Considerar la falsificación o piratería corno delito penai si se realiza a escala comercial • Las sanciones deben ser suficientemente severas/ disuasorias. Periodos transitorios de aplicación. • 1 año 1 año para los países desarrollados (hasta el I de enero de 1996). • 5 años para los PEO (hasta el I de enero de 2000). ,c:x • 1 1 af\os para los PMA inicialmente (hasta el 1 de enero de 2006, y hasta 2016 para el caso de las patentes de medicamentos). Este periodo se ha ampliado posteriormente hasta julio de 2034. y I de enero de 2033 en el caso de patentes para medicamentes (de confonnidad con la Declaración de Doha relativa al Acuerdo sobre los ADPlC y la Salud Pública de 2001). Los países que se unan a la OMC deben cumplir con el TRIPS desde el momento de su adhesión, no se benefician de un periodo transitorio.




\clearpage
\section{Organización Mundial del Comercio: Acuerdos Plurilaterales no relativos a mercancias}
\puntosclave{
    Ante el estancamiento de la Ronda de Doha y de las negociaciones multilaterales, los Miembros han empleado hasta tres estrategias distintas para avanzar en temas distintos de los de mercancías: acuerdos de cosecha temprana en la OMC (que han sido escasos por el     momento), y fuera de la OMC mediante negociaciones bilaterales (con un país o grupos de países) y las negociaciones plurilaterales (nos centramos en las últimas por su potencial de multilateralización)
    
    Las vías mas prometedoras de avance actualmente son las siete Iniciativas Conjuntas Plurilaterales ($4$ lanzadas en la CM de Buenos Aires, y las tres medioambientales lanzadas entre 2020 y 2021 ). De ellas, las negociaciones de la ICP de Reglamentación Nacional ya están cerradas, y solo queda trasladar dichos compromisos al GA Ts para su aplicación .
}

Como hemos podido apreciar, desde la creación de la OMC el número de temas de negociación ha sido muy extenso. Sin embargo, a pesar de haberse creado numerosos comités y grupos de trabajo de temas diversos la mayor parte de las negociaciones han prosperad\footnote{
Por ejemplo comercio y medio ambiente, los cuatro temas Singapur (comercio y competencia, comercio e inversiones, contratación pública, y facilitación de comercio) y/o comercio electrónico. Sobresale a estos fracasos el acuerdo de Facilitación de Comercio, acuerdo multilateral sobre mercancías que se negoció en 2012 y entró envigor en 2017.
} 

\subsection{Acuerdos Plurilaterales: Los frutos del \textit{single undertaking}}
Frente a este estancamiento, la OMC ha ido cambiando su enfoque de negociación pasando a adoptar el \textit{Single Undertaking}, iniciado en la Ronda de Doha. Este enfoque más prágmatico, pero más controvertido, es un inicio de exploración por la vía plurilateral cuyos resultados vamos a presentar ahora. 

En diciembre de 2011, durante la CM$8$ en Ginebra, se reconoció el estancamiento de la Ronda de Doha y se dio una orientación política para sustituir el objetivo de compromiso único de las negociaciones (cerrar la Ronda con un gran acuerdo en todos los temas incluidos en la Ronda) por otro más realista de llegar a \textbf{acuerdos provisionales o definitivos basados en el consenso antes de la conclusión total del compromiso único}. 


\subsubsection{¿Cuándo ha funcionado? En vigor}
Los frutos más importantes de esta nueva vía son tres acuerdos plurilaterales, dos relativas al comercio de mercancías y uno relativo a la contratación pública.

\paragraph{Acuerdo de Contratación Pública (ACP, 1994; 2012)}
El Acuerdo de Contratación Pública (ACP) es un acuerdo negociado en el marco de la OMC, en la Ronda de Doha, que adoptó la negociación y adopción plurilateral por ser su materia objeto de excepción tanto en el GATT como en el GATS. Especialmente importante tanto por sus efectos cualitativos sobre la eficiencia del sector público como por la importancia cuantitativa que el mercado de contratación pública representa. Actualmente cuenta con $21$ partes, o $48$ miembros de la OMC, ya que la UE es parte, e incluye a $27$ miembros de la OMC más la UE, de ahí la discrepancia en los números. Sus antecedentes directos se remontan al \textit{Código de la Ronda de Tokio sobre Compras del Sector Público}, adoptado en la Ronda de Tokio (1979), pero con la creación de la OMC se incorporó el nuevo Acuerdo de Contratación Pública (1994) revisado y modificado posteriormente en 2012, que entró en vigor hasta 2014.

Su objetivo fundamental es la \textbf{apertura de mercados de contratación pública entre partes}, así como la mejora de la transparencia y la competencia en este. Como resultado de varias rondas de negociaciones el volumen de actividad de contratación entre partes asciende a un vallor estimado $1,7\text{B}\$$ anuales. Consta de dos partes:
\begin{la}
    \item \textbf{Texto del Acuerdo}. El texto del Acuerdo establece normas de garantía de la competencia, el trato equitativo y la transparencia en la esfera de la contratación pública. Entre otros, establece el prinicpio de \textbf{publicidad y transparencia} y el principio de \textbf{no discriminación}. No obstante, esas normas no se aplican automáticamente a todas las actividades de contratación de cada Parte;
    \item \textbf{Listas de compromisos de las partes} (en materia de Acceso al Mercado). Las listas de cobertura desempeñan un papel decisivo en la determinación de si una actividad de contratación está abarcada o no por el Acuerdo. En esencia, las Partes se comprometen a identificar las entidades que llevan a cabo actividades de contratación y publicar aquellas licitaciones de bienes, servicios o servicios de construcción que superen un determinado umbral.
\end{la}

El tratado está administrado por el Comité de Contratación Pública, integrado por representantes de todas las partes y existen dos mecanismos para hacer cumplir el Acuerdo: el mecanismo de recurso interno a nivel nacional y el mecanismo de solución de diferencias de la OMC a nivel internacional.


\paragraph{Acuerdo sobre Facilitación de las Inversiones para el Desarrollo (AFID)}


\subsubsection{¿Cuándo no ha funcionado? Negociaciones infrucuosas}
Existen otros dos acuerdos importantes fruto del marco plurilateral, pero cuyas negociaciones están actualmente paradas. Uno en el ámbito del comercio de mercancias y otro que surgió con la idea de suceder al GATS.\footnote{El relativo a mercancias es el Acuerdo EGA (Environmental Goods Agreement) sobre bienes medioambientales. Sus negociaciones se iniciaron en 2014 y están actualmente paradas. Están presididas por Australia, con la participación de Canadá, China, Taipei Chino, Costa Rica, la Unión Europea y sus $28$ estados miembros, Hong Kong, Islandia, Israel, Japón, República de Corea, Nueva Zelanda, Noruega, Singapur, Suiza, Turquía y los Estados Unidos.

Su objetivo es reducir los aranceles sobre \textbf{bienes verdes} que benefician al medio ambiente para reducir costes y beneficiar el comercio. El acuerdo tiene como objetivo ampliar la lista de $54$ bienes ambientales que los líderes de APEC acordaron para la reducción de aranceles en 2012. 
} 

\paragraph{Trade in Services Agreement: ¿Sucesión del GATS?}
El acuerdo TISA (\textit{Trade in Services Agreement}) de comercio de servicios es un acuerdo plurilateral para la liberalización de los servicios en el que participan $23$ miembros de la OMC, incluida la UE. Las negociaciones comenzaron en marzo de 2013 y están actualmente paradas (la última ronda fue en diciembre de 2016), aunque no se han suspendido oficialmente. 

Después de la entrada en vigor del GATS, los miembros de la OMC continuaron las negociaciones en algunos sectores, que permitieron acordar disciplinas adicionales posteriores al cierre de la R. Uruguay: Servicios Financieros (1997), telecomunicaciones (1998) y servicios profesionales de contabilidad (1998). Estos compromisos posteriores se integraron al GATS.

Una vez que se lanzó la Ronda de Doha, las negociaciones de servicios se integraron en las negociaciones globales de la Agenda de Doha en 2001. El formato de negociación bilateral con el sistema de ofertas y demandas fue el más utilizado al principio de la ronda, aunque también se intentaron otros formatos plurilaterales y multilaterales. No obstante, el desacuerdo entre los países con interés en liberalizar el acceso a mercado agrícola (PEO) y NAMA (PO) centró la atención de las negociaciones y los servicios fueron prácticamente olvidados. 

En julio de 2008, durante la Conferencia mini-ministerial de la OMC en Ginebra, los países estuvieron cerca del cierre de Doha. Finalmente, a pesar de los progresos realizados, el cierre de la Ronda fracasó debido a desacuerdos sobre las normas agrícolas y el acceso a los mercados, estando las negociaciones estancadas desde entonces. 

De esta idea nacieron las negociaciones de TISA. El acuerdo TISA El Acuerdo de Comercio de Servicios (TiSA- Trade in Services Agreement) es un acuerdo comercial plurilateral que, ante el estancamiento de la Ronda de Doha, trató de avanzar en la liberalización de los servicios. Las negociaciones comenzaron en marzo de 2013. Incluye a 23 miembros de la OMC, que suponen un 70\% del comercio mundial de servicios.\footnote{Australia, Canadá, Chile, Taiwán, Colombia, Costa Rica, la UE, Hong Kong, Islandia, Israel, Japón, Corea del Sur, Liechtenstein, Mauricio, México, Nueva Zelanda, Noruega, Pakistán, Panamá, Perú, Suiza, Turquía y los EEUU. China no forma parte de estas negociaciones
}

• Está diseñado para ser compatible con la OMC, basado en GATS • Es un acuerdo abierto a la incorporación de nuevos participantes y multilateralizable: una vez  se alcance una masa critica de participantes, se podría multilateralizar incorporando los compromisos alcanzados al GATS. 

Es un acuerdo amplio, cubre todos los servicios (igual que GATS, con sus mismas exclusiones) • No se aplica NMF ya que se diseñó como acuerdo de integración económica, al ser un acuerdo de liberalización y no sectorial (incluye casi todos los sectores de servicios, al igual que GATS). 

Utilizar un formato mixto, de lista positiva para Acceso a Mercado y lista negativa para Trato Nacional 

Desde la última ronda de negociaciones de diciembre de 2016 no se han producido novedades significativas.










\subsection{Iniciativas Plurilaterales}
Paralelamente, a raíz de la $11ª$ Cumbre Ministerial de Buenos Aires 2017, se han lanzado lo que se conoce como \textbf{Iniciativas Conjuntas Plurilaterales} (ICP). Estas versan, fundamentalmente, sobre cuatro materias: el comercio electrónico ($87$ miembros que suponen el $90\%$ del comercio mundial), las MiPYMEs (microempresas, y pequeñas y medianas empresa), la facilitación de inversiones para el desarrollo y la reglamentación nacional de servicios.\footnote{La iniciativa para la \textbf{facilitación de inversiones para el desarrollo} cuenta con más de 110 participantes y pretende mejorar el entorno empresarial y de inversión. Se abordan cuestiones de transparencia, requisitos de información, y agilización de procedimientos administrativos aplicables a la inversión extranjera y luchas contra la corrupción. No incluye negociaciones de acceso a mercado ni protección de inversiones o resolución de controversias. Muchas de las disposiciones se están redactando en forma de \textbf{mejores esfuerzos}. Es decir, como no vinculantes, un nivel de ambición limitado que era de esperar ya que, al aplicar la clausula NMF, lo que los miembros acuerden en este acuerdo plurilateral se extenderá al resto de miembros no participantes. Las negociaciones comenzaron formalmente en 2020, y se busca finalizar las negociaciones para finales de 2022.

Por otro lado, en la \textbf{iniciativa de comercio electrónico} se debaten sobre seis temas principales: posibilitación del comercio electrónico, apertura y comercio digital, confianza y comercio digital, cuestiones transversales, telecomunicaciones y acceso a mercado.
}

Las iniciativas plurilaterales participan sólo algunos miembros de la OMC, pero permanecen abiertas a cualquier miembro que desee unirse y permiten avanzar en temas cuyas negociaciones están bloqueadas en el ámbito multilateral. 

\paragraph{Reglamentación Nacional en Servicios}
La iniciatia plurilateral para la reglamentación nacional de servicios es, de momento, la primera en la que se ha llegado a un acuerdo. Cuenta con $70$ participantes (incluida la UE) y se inició en base al mandato de negociación incluido en GATS, por lo que comparte el mismo objetivo que el grupo de trabajo de reglamentación nacional de la OMC. 

Las negociaciones se cerraron en septiembre de 2021 y quedaron manifiestas en el \textbf{Documento de Referencia de Reglamentación Nacional} de Servicios. Incluye acuerdos sobre transparencia y mejoras en los procedimientos y requisitos para la concesión de licencias, en todos aquellos sectores en los que los miembros han realizado concesiones de liberalización bajo el GATS. 

\paragraph{Iniciativas ambientales}
Posteriormente entre 2020 y 2021 se han lanzados nuevas ICP medioambientales sobre tres temas en los que se está trabajando en la actualidad:
\begin{la}
    \item Debates Estructurados sobre \textbf{Comercio y Sostenibilidad Ambiental}. Lanzada en noviembre de 2020, actualmente está liderada por Canadá y Costa Rica e incluye a $74$ Miembros que representan el $80\%$ del comercio mundial. En la CM$12$ se ha informado de los progresos efectuados en los cuatro grupos de trabajo informales establecidos: bienes y serv1cios medioambientales, medidas sobre el clima relacionadas con el comercio, economía circular, y subvenciones;\footnote{Podría sustituir o terminar las negociaciones iniciadas con el Acuerdo de bienes medioambientales (EGA), pero ampliando a servicios (EGS-Environmental Goods and Services).
    }
    \item Diálogo Informal sobre \textbf{Contaminación por Plásticos y Comercio de Plásticos Ambientalmente Sostenible}. Lanzada en noviembre de 2020, el plan de trabajo se aprobó en 2022. Busca reducir la contaminación de plásticos y conseguir un modelo de comercio de plásticos más sostenible. Lideran la iniciativa Australia, Barbados, China, Ecuador, Fiji, y Marruecos, y forman parte de ésta $75$ miembros que representan el $75\%$ del comercio mundial de plásticos; y,
    \item \textbf{Reforma de los Combustibles Fósiles}. Lanzada en diciembre de 2021, está liderada por Nueva Zelanda. Su objetivo es racionalizar y eliminar progresivamente las subvenciones ineficientes a los combustibles fósiles. Cuenta con $47$ Miembros.
\end{la}





























\clearpage
\section{Conclusión} 


\subsection{Recapitulación}


\subsection{Extensiones y relación con otras partes del Temario}



\subsection{Opinión}

\subsubsection{Idea final (Salida o cierra)}

\clearpage
\pagenumbering{gobble}
\MyBibliography



\MyBibItem{DAWID y GATTI}{Chapter 2 - Agent-Based Macroeconomics}{2018}{https://doi.org/10.1016/bs.hescom.2018.02.006}


% ANEXOS
\clearpage
\anexo{\textbf{Preguntas Test}}
%\input{\TemaCode/questions.tex} 



\end{document}